% Firmware Dump Emulation Workflow
%
% Shows the 5-step pipeline for emulating a firmware .bin dump:
% SVD Parsing -> Peripheral Set -> Board YAML -> QEMU Run -> MMIO Report
%
% Compile with: pdflatex firmware_dump_workflow.tex
% Convert to PNG: pdftoppm -png -r 200 firmware_dump_workflow.pdf firmware_dump_workflow
%
% Author: Mathieu Renard <mathieu.renard@twistedwires.io>
% Copyright (C) 2026 Twisted Wires Security Lab
% SPDX-License-Identifier: Apache-2.0

\documentclass[tikz,border=10pt]{standalone}
\usepackage{tikz}
\usetikzlibrary{arrows.meta,positioning,shapes.geometric,fit,backgrounds,calc}

\begin{document}

\definecolor{trustzoneblue}{RGB}{0,123,255}
\definecolor{securegreen}{RGB}{34,139,34}
\definecolor{spiblue}{RGB}{0,102,204}
\definecolor{i2corange}{RGB}{230,126,34}
\definecolor{uartgreen}{RGB}{39,174,96}
\definecolor{gpiored}{RGB}{192,57,43}
\definecolor{mcugray}{RGB}{100,100,100}
\definecolor{passgreen}{RGB}{40,167,69}

\begin{tikzpicture}[
    step/.style={
        rectangle, rounded corners=4pt, draw=#1, fill=#1!10,
        text width=3.2cm, minimum height=2.2cm, align=center,
        font=\sffamily\small, line width=1pt
    },
    input/.style={
        rectangle, rounded corners=2pt, draw=mcugray, fill=mcugray!5,
        text width=2.8cm, minimum height=0.8cm, align=center,
        font=\sffamily\tiny, line width=0.5pt
    },
    arrow/.style={
        -{Stealth[length=5pt]}, line width=1.2pt, color=#1
    },
    bigarrow/.style={
        -{Stealth[length=8pt]}, line width=2pt, color=mcugray!60
    },
]

% Step 1: Firmware + SVD
\node[step=spiblue] (step1) {
    \textbf{1. Identify Target}\\[4pt]
    \texttt{firmware.bin}\\
    \texttt{STM32F405.svd}\\[2pt]
    {\tiny Vector table analysis}
};

% Step 2: Peripheral Set
\node[step=i2corange, right=1.2cm of step1] (step2) {
    \textbf{2. Peripheral Set}\\[4pt]
    \texttt{MCU\_REGISTRY}\\[2pt]
    {\tiny 21 supported MCUs}\\
    {\tiny or SVD auto-stub}
};

% Step 3: Board YAML
\node[step=securegreen, right=1.2cm of step2] (step3) {
    \textbf{3. Board Config}\\[4pt]
    \texttt{board.yaml}\\[2pt]
    {\tiny MCU + clock + SRAM}\\
    {\tiny + external devices}
};

% Step 4: QEMU Emulation
\node[step=trustzoneblue, right=1.2cm of step3] (step4) {
    \textbf{4. QEMU Run}\\[4pt]
    \texttt{qemu-system-arm}\\
    \texttt{+ TCP server}\\[2pt]
    {\tiny Firmware executes}
};

% Step 5: MMIO Report
\node[step=uartgreen, right=1.2cm of step4] (step5) {
    \textbf{5. Analysis}\\[4pt]
    MMIO Trace\\
    UART Output\\
    LaTeX Report\\[2pt]
    {\tiny + GDB debug}
};

% Arrows between steps
\draw[bigarrow] (step1) -- (step2);
\draw[bigarrow] (step2) -- (step3);
\draw[bigarrow] (step3) -- (step4);
\draw[bigarrow] (step4) -- (step5);

% Input files below step 1
\node[input, below=0.6cm of step1] (fw) {
    \texttt{.bin} firmware dump\\
    {\tiny Initial SP + Reset\_Handler}
};
\draw[arrow=spiblue] (fw) -- (step1);

% SVD file below step 2
\node[input, below=0.6cm of step2] (svd) {
    \texttt{.svd} register defs\\
    {\tiny Peripheral base addresses}
};
\draw[arrow=i2corange] (svd) -- (step2);

% YAML below step 3
\node[input, below=0.6cm of step3] (yaml) {
    \texttt{name: MyBoard}\\
    \texttt{mcu: STM32F405}
};
\draw[arrow=securegreen] (yaml) -- (step3);

% GDB below step 4
\node[input, below=0.6cm of step4] (gdb) {
    GDB on port 1234\\
    {\tiny Fault handler breakpoints}
};
\draw[arrow=trustzoneblue, dashed] (gdb) -- (step4);

% Output files below step 5
\node[input, below=0.6cm of step5] (out) {
    \texttt{trace.mmio}\\
    \texttt{report.tex (TikZ)}
};
\draw[arrow=uartgreen] (step5) -- (out);

% Title
\node[above=0.8cm of step3, font=\sffamily\large\bfseries, color=mcugray] {
    Firmware Dump Emulation Workflow
};

\end{tikzpicture}
\end{document}
